\subtituloA{Intérprete de JavaScript y compilación JIT}

definicion{Intérprete de JavaScript}{Módulo que analiza, compila con técnicas JIT (Just-In-Time) y ejecuta código para añadir dinamismo y comportamiento interactivo a la página web.}
subtituloA{Peticiones asíncronas con fetch}
begin{codigo}
fetch(\"api/datos.json\")
  .then(function (respuesta) {
    if (!respuesta.ok) throw new Error(\"Respuesta no satisfactoria: \" + respuesta.status);
    return respuesta.json();
  })
  .then(function (datos) {
    console.log(\"Datos obtenidos del servidor:\", datos);
  })
  .catch(function (error) {
    console.error(\"Error en la comunicación de red:\", error.message);
  });




 nota{Para el examen}{ En el código fetch anterior, ¿qué ocurre si la \textbf{Respuesta:} del servidor no es ok y cómo se maneja el error? }

"

  ]

Wait, I need to check the rules again carefully.

Rule 3: "El primer archivo de cada subtema debe incluir  y 

...

. Los archivos siguientes usan ." In my draft, the first file has \cmd{\secciongrande} and \cmd{\resumen}, good. The second and third have \cmd{\subtituloA}, good.

Rule 10: "SIN LÍNEAS EN BLANCO DENTRO DE MACROS: Dentro de 

...

{
... " I need to ensure no empty lines inside these macros. In my LaTeX strings, I have newlines outside the macros? Actually, the \cmd{latex} field content is a string. I need to write it such that when compiled, there are no empty lines inside the macros.

The way LaTeX works in this context, the \cmd{}

...

 macro content should not contain blank lines. I'll write the content as a single line or with careful line breaks, but since it's a JSON string, I can include \cmd{}
} but must avoid blank lines within the macro arguments. The rule says "NO dejes líneas en blanco dentro de \

...

". So inside the braces, no empty lines. I'll make sure the text inside \cmd{}

...

 is continuous without empty lines. Similarly for \cmd{\nota{...}{
...
}}

\cmd{\definicion{...}{...}}

\cmd{\lista{}
  ...
}.

Rule 8: "CÓDIGO ÍNTEGRO Y SIN DIVIDIR: Todo bloque de código en cualquier lenguaje debe colocarse completo en un único bloque

\begin{codigo} ...
\end{codigo}



NUNCA fragmentado o dividido con texto explicativo intercalado. " I have \cmd{}

\begin{codigo}...
\end{codigo}



 outside the macros, good.

I need to ensure the code is complete and not inside other macros.

Rule 7: "VIÑETAS Y DIVISIÓN OBLIGATORIA (PROHIBIDO PEGAR ELEMENTOS Y PROHIBIDO ): En \lista{
  **Elemento 1:** ... |
  **Elemento 2:** ...
}.  with the \cmd{|} separator format. I need to make sure I don't use \cmd{}.

.

Rule 6: "DEFINICIONES OBLIGATORIAS: Aporta SIEMPRE definiciones formales con para cada concepto fundamental del material." I have \cmd{\definicion} in each file, good.

Rule 9: "PREGUNTAS DE AUTOEVALUACIÓN Y EXAMEN DEL PROPIO APARTADO: Cada caja \nota{Para el examen}{
...
} (o > **Para el examen:**) DEBE formular preguntas que pertenezcan ESTRICTAMENTE A LOS CONCEPTOS TRATADOS EN ESTE APARTADO CONCRETO.

Escribe siempre \nota{Para el examen}{
Pregunta... " I have \cmd{\nota{Para el examen}{
... }} in each, good.

Rule 5: "ESTÁ TERMINANTEMENTE PROHIBIDO usar fórmulas de relleno vacías o metatexto como «Este apartado explora cómo...», «En este apartado se describe...», «Este apartado analiza...», «Además se incluyen ejemplos...». Enseña DIRECTAMENTE la teoría y los hechos técnicos desde la primera línea." I'll avoid such phrases.

Rule 4: "PROHIBIDO inventar subtítulos genéricos o de relleno («Elementos principales», «Componentes principales», «Aspectos clave», «Idea clave de cada apartado», «Puntos clave», «Conceptos principales»...). Un subtítulo debe nombrar el concepto real del material." I'll make sure subtitles are real concepts.

Rule 2: "Todo el contenido debe estar 100% en ESPAÑOL ESTÁNDAR (castellano). Prohibido utilizar palabras en catalán en títulos, slugs o explicaciones." I'll translate everything to Spanish. The PDF name is Catalan, but I'll use Spanish terms. E.g., "interfaz de usuario", "motor de renderizado", "intérprete de JavaScript", etc.

Rule about \cmd{titulo_subtema}: The top-level \cmd{titulo_subtema} should be "Título formal del apartado en español" but I need to put a real title. The template shows it as a placeholder, but the rules say it's part of the JSON format and must be filled. I'll set it to "Arquitecturas y herramientas en entorno cliente" or similar.

Rule about \cmd{tema_existente}: The template has \cmd{0}. I'll keep it \cmd{0}.

Now, let's refine the LaTeX content to strictly follow the no-empty-lines rule.

For \cmd{}

...
}: The content must be 3-5 lines of substantial theory, no empty phrases. I'll write a concise summary.

For \cmd{\definicion{...}{...}}: Formal definition.

For \cmd{}: Real concept title.

: Items separated by \cmd{ | }, no empty lines inside the braces. The rule says: "En \lista{
  **Elemento 1:** ... |
  **Elemento 2:** ... " So the format is exactly that, with items and \cmd{|} between them. I'll make sure there are no line breaks inside the braces that create empty lines. I can have the items on separate lines if I use \cmd{\\} but that might create empty lines. Better to keep them in one line or use the \cmd{|} separator on one line. The rule says "NO dejes líneas en blanco dentro de ...
}". So as long as there are no completely empty lines, it's fine. I'll format as \cmd{\lista{}
  **Elemento 1:** Explicación técnica. |
  **Elemento 2:** Explicación técnica.

}} all on one line within the macro.

For \cmd{\nota{Para el examen}{
...
}}: Question and reasoned answer.

. I'll include it, but need to check where to place it. The system prompt example has it as a standalone macro, not inside another. I'll place it at the end of the \cmd{latex} content, or maybe outside the main flow?

The rules don't prohibit it, and the system prompt example includes it. I'll add it after the \cmd{\nota} or within the file's latex, but ensuring it's properly formatted. I'll put it as a separate macro after the \cmd{\nota}, but within the \cmd{latex} string.

However, the \cmd{latex} field is a single string; I can include multiple macros. .

Now, let's think about the \cmd{conceptos_clave} array. I'll put 2 key concepts.

Also, the \cmd{nombre} field: "slug corto de 1 a 3 palabras clave conceptuales en español estándar, ej: "servidor", "modelos-ejecucion-codigo", "tecnologias-entornos-cliente", "navegador-web", "capacidades-limitaciones", "herramientas-cliente"; NUNCA slugs en inglés como "models-execution", "client-server" ni términos en catalán como "model", "entorn", "eines", "clients", "execucio", "fitxers"". I'll create Spanish slugs like "arquitectura-navegador", "motor-renderizacion", "intprete-javascript".

The \cmd{subtema} field: "1" for all files since we're doing subtema 1.

The \cmd{titulo} field: "Título formal del apartado en español" - I'll put the actual title of the section/file.

Now, let's draft the JSON carefully, ensuring every macro is closed with braces, no empty lines inside, etc.

I'll start with the JSON skeleton:

```json
{
  "tema\_existente": 0,
  "titulo\_tema": "Tema 1",
  "titulo\_subtema": "Arquitecturas y herramientas en entorno cliente",
  "conceptos\_clave": ["Arquitectura de navegador web", "Motor de renderizado y ejecución de código"],
  "archivos}
